\section{Background}
\label{sec:Background}

\subsection{Zero-Knowledge Proofs}
Today, the state-of-the-art ZKP protocols are zero-knowledge Succinct Non-interactive Arguments of Knowledge (zkSNARKs). They have three key properties:
\begin{itemize}
    \item {\textbf{zero-knowledge:}} The verifier learns nothing beyond the validity of the statement.
    \item {\textbf{succinctness:}} The proof is short and verified quickly.
    \item {\textbf{non-interactivity:}} The proof is a single message, eliminating the need for back-and-forth communication.
\end{itemize}

Modern zkSNARKs combine an Interactive Oracle Proof (IOP) and a Polynomial Commitment Scheme (PCS). The IOP encodes the witness and computation as polynomials, which the verifier checks via probabilistic queries. Applying the Fiat-Shamir heuristic~\cite{fiat_shamir} makes public-coin IOPs non-interactive. The PCS lets the prover commit to a polynomial and later open it at chosen points, ensuring both binding (the prover cannot change the committed polynomial) and hiding (nothing about the polynomial is revealed until it is opened).

Different zkSNARKs balance proving time, proof size, and security assumptions based on their IOP and PCS designs. Groth16~\cite{groth16} offers fast verification and tiny proofs, ideal for blockchains, but requires a circuit-specific (i.e., per-application) trusted setup and scales poorly with circuit size. Code-based protocols~\cite{orion} use error-correcting codes and Merkle commitments to achieve transparency and post-quantum security, but at the cost of multi-megabyte proofs and slower verification. HyperPlonk~\cite{hyperplonk} trades off slightly larger proofs for universal setup and linear-time proving.


\subsection{Multi-Scalar Multiplication}
Multi-Scalar Multiplication (MSM) is a key kernel used in \textit{polynomial commitments} and relies on elliptic curve cryptography.
MSMs involve computing a linear combination (i.e., dot product) of scalar and (2D or 3D) point pairs on an elliptic curve.
Given scalars \( k_1, k_2, \ldots, k_n \) and corresponding elliptic curve points \( P_1, P_2, \ldots, P_n \), the goal is to efficiently compute $s = \sum_i k_iP_i$.
MSMs are computationally intensive because each term \( k_iP_i \) involves a point-scalar multiplication -- a costly elliptic curve operation. 
MSMs constitute a significant portion of the total runtime in many zkSNARKs, including Groth16 and HyperPlonk. 
As a result, prior work has focused on accelerating MSMs using both application-specific and programmable platforms~\cite{szkp, pipezk, cuzk, gzkp, distmsm, priormsm, rezk}.
Rather than computing each scalar multiplication directly, optimized approaches such as Pippenger’s algorithm~\cite{pippenger} are commonly employed.
Pippenger’s algorithm improves efficiency by restructuring the computation to rely more heavily on faster operations like point addition and point doubling.

\subsection{SumCheck}
\label{sec:sumcheck_intro}
SumCheck~\cite{sumcheck} is a round-based interactive protocol where a prover $\mathcal{P}$ convinces a verifier $\mathcal{V}$ that the sum of a multivariate polynomial, $f(X_1, \ldots,X_\mu)$, over the \textit{boolean hypercube} (all boolean values of $X_i \ldots X_\mu$) equals a claimed value $C$.
The protocol proceeds in $\mu$ rounds, one per variable. 
In round $i$, $\mathcal{P}$ sends a univariate polynomial
\begin{equation*}
    s_i(X_i) = \sum_{x_{i+1} \in \{0,1\}} \cdots \sum_{X_\mu \in \{0,1\}} 
    f(r_1, \ldots, r_{i-1}, X_i, \ldots, X_\mu)
\end{equation*}
where $r_1, \ldots, r_{i-1}$ are $\mathcal{V}$’s random challenges from prior rounds. $\mathcal{V}$ checks that $s_i(0) + s_i(1)$ equals the previous round's claim (or the initial claim $C$ for round $1$).
$\mathcal{V}$ then chooses a random field element $r_i$ and sends it to $\mathcal{P}$, who sets $X_i = r_i$ for the next round.
After $\mu$ rounds, $\mathcal{V}$ evaluates $f$ at $(r_1, \ldots, r_\mu)$ and accepts if all checks pass.

ZKPs execute SumCheck over \textit{multilinear extensions} (MLEs) -- polynomials with degree at most $1$ in each variable. MLEs represent multivariate functions compactly and, in hardware, can be stored as flat lookup tables indexed by binary inputs. Throughout this paper, we use the terms ``MLEs," ``MLE tables," and ``polynomials" interchangeably.

\subsubsection{\textbf{Polynomials and Gates}}
In ZKPs, polynomials encode and help enforce the correctness of the computation being proven. A popular encoding scheme is Plonk \cite{gabizon2019plonk}. Plonk's \textit{arithmetization} scheme converts the underlying computation into an arithmetic circuit consisting of gates. Plonk circuits support additions, multiplications, conditionality, and equality checks. They can be represented by the composite polynomial 
\begin{equation*}
\label{eq:plonk}
    f_{plonk} = q_Lw_1 + q_Rw_2 + q_Mw_1w_2 - q_Ow_3 + q_C
\end{equation*}
The composite polynomial $f_{plonk}$ has 5 terms and a total of 8 unique, constituent polynomials.
These constituent polynomials are the selectors (or enables) $q_i$ and witnesses $w_i$, and they are all multilinear.
Notably, 
% $f$
$f_{plonk}$
is not itself multilinear; rather, it is composed of multilinear polynomials.

Each of these polynomials encodes a function of a multivariate input vector $X=[X_1, X_2, \ldots, X_\mu]$, which encodes the gate index in binary. For example, if $\mu = 3$, there are 8 gates, and gate index 3 would be represented by evaluating $f_{plonk}(1, 1, 0)$. For a multiply operation, we would set $q_M(1, 1, 0) = 1$, $q_O(1, 1, 0) = 1$ and the other enable polynomials to $0$. Then, gate 3 would execute correctly if and only if $f_{plonk}(1, 1, 0) = 0$, that is, 

\begin{equation*}
w_1(1,1,0)\cdot w_2(1,1,0) - w_3(1,1,0) = 0
\end{equation*}
Representing circuits as polynomials lends itself naturally to SumCheck-based provers where we can check (with additional constraints) that each gate $f_{plonk}(X)$ evaluates to $0$ and that
\begin{equation*}
\sum_X f_{plonk}(X) = 0 
\end{equation*}

\subsubsection{\textbf{High-degree, Expressive Gates}}
The challenge of using Plonk gates is that all computations must be mapped to only what is supported by the Plonk polynomial, 
similar to mapping program code to a RISC ISA.
In some cases, this results in circuits that require millions ($>2^{20}$) of gates, and in turn, polynomials with millions of evaluations.
Recent protocols and libraries ~\cite{halo2, hyperplonk} support \textit{expressive, high-degree} gates, which encode complex operations and significantly reduce gate count.
For example, HyperPlonk’s \textit{Jellyfish} gate captures common patterns such as hashing; across benchmarks, it reduces gate counts by up to 32× compared to less expressive \textit{Vanilla} Plonk gates.
Similarly, Halo2 uses custom gates for elliptic curve operations. However, these expressive gates increase the complexity of the SumCheck algorithm used to verify them.



\subsubsection{\textbf{SumCheck on Multilinear Polynomial Products}}
SumCheck over a single multilinear polynomial is relatively straightforward.
Consider a polynomial $f(X_1, X_2, X_3)$ with $\mu = 3$ and its MLE table representation.
In the first round, the prover computes the univariate polynomial $s_1(X_1)$ by summing entries with $X_1 = 0$ and $X_1 = 1$, yielding two evaluations $s_1(0)$ and $s_1(1)$ sent to the verifier.
The prover then computes an \textit{MLE Update} by fixing $X_1$ to a random challenge $r_1$, typically generated by hashing the round evaluations (e.g., with SHA3).
For each $(X_2, X_3)$, the entries $f(0, X_2, X_3)$ and $f(1, X_2, X_3)$ define a line
in $X_1$, which is extrapolated to $r_1$: $f(r_1,X_2, X_3) = f(0, X_2, X_3)\cdot(1-r_1) + f(1, X_2, X_3)\cdot r_1$.
This halves each MLE table before the next SumCheck round.

In recent SumCheck-based ZKP protocols \cite{spartan, orion, hyperplonk, jolt}, we typically work with multi-term compositions of multilinear polynomials, like $f_{plonk}$ in \autoref{eq:plonk}. In these protocols, we are given \textit{only} the constituent polynomials and their composition structure, and then must perform SumCheck over the composition. These SumChecks are naturally more complicated.
In general, to run SumCheck on a $d$-degree polynomial, we require $d+1$ evaluations for $s_i(X_i)$ after each $i^{th}$ round of SumCheck.
For example, $f_{plonk}$'s overall degree is $3$, so \textit{$4$} evaluations are needed per SumCheck round.
These evaluations must be computed for each individual MLE before computing term-wise products and summing across terms.

\autoref{fig:SumCheck_dataflow_example} shows the dataflow for one polynomial product term.
Each pair of entries in an MLE table actually represents the evaluations at $X_i = 0$ and $1$; these values are then \textit{extended} to $X_i = 2, \ldots, d$, using the same formula as MLE Update (hence multilinear \textit{extension}).
Then, the extensions for all polynomials in a term are multiplied together.
These products are summed down the MLE table, yielding $d+1$ values for a term.
When there is more than one term, 
this process will be repeated for each term, with the $d+1$ evaluations summed across terms to form the final polynomial $s_i(X_i)$ for round $i$.

We now present our generalizable SumCheck accelerator.




